\section{Related work}\label{related-work}

\subsection{Grinding roughness monitoring with AE and vibration}\label{rw:ae-vibration}

Grinding differs from many cutting processes because material removal is distributed over a large number of abrasive grains with stochastic protrusion heights, local rubbing, ploughing, chip formation, grain wear, and possible thermal damage. Surface roughness is therefore affected by controllable process parameters \cite{malkin2008grinding} and evolving wheel--workpiece interaction mechanisms \cite{Brinksmeier2006}. Reviews of surface integrity in material-removal processes have emphasized that topography, residual stress, microstructural alteration, and thermal-mechanical effects jointly determine functional performance, which makes roughness prediction a useful but incomplete indicator of final quality \cite{Jawahir2011}. In-process sensing is consequently attractive because it can provide information on process evolution before offline metrology is available.

AE has been used in machining monitoring because high-frequency stress waves are generated by localized contact \cite{Dornfeld1992}, micro-fracture, and rubbing interactions \cite{Teti2010}. In grinding, AE is often sensitive to events near the wheel--workpiece interface and can respond to wheel loading, dressing condition, contact initiation, burn-related changes, and changes in removal regime. However, AE signals are filtered by the sensor, couplant, mounting path, fixture, and machine structure. Interpretation of frequency-localized attributions therefore requires caution, especially when model explanations emphasize regions above the nominal flat-response band of the sensor.

Vibration sensing provides a complementary route because it measures the dynamic response of the machine-tool structure rather than only localized high-frequency activity. Chatter, spindle modes, and machine compliance can leave signatures in vibration spectra \cite{Altintas2004}, which can also reveal wheel wear or imbalance \cite{Quintana2011}. Empirical grinding studies have shown that vibration features can correlate with surface roughness in cylindrical grinding, supporting the use of vibration as a roughness-prediction modality \cite{Hassui2003}. Because AE and vibration emphasize different frequency ranges and physical pathways, their fusion can reduce reliance on a single sensor and can expose disagreements between local contact activity and system-level dynamics.

\subsection{Signal representations for grinding monitoring}\label{rw:representations}

Feature engineering in grinding and machining monitoring has typically combined time-domain, frequency-domain, and time--frequency descriptors. Time-domain features such as RMS, variance, skewness, kurtosis, peak-to-peak value, crest factor, and envelope statistics are compact and inexpensive. Frequency-domain features such as band power, dominant-frequency amplitude, spectral centroid, and spectral kurtosis can isolate energy changes associated with dynamic modes or impulsive events \cite{Antoni2006}. Wavelet-based features have also been used because grinding and AE signals are non-stationary and may contain transient events that are not well represented by pass-level averages alone \cite{Teti2010}.

Spectrograms provide a direct two-dimensional representation of temporal and spectral variation. In deep-learning pipelines, they allow CNNs to learn local time--frequency patterns using architectures originally developed for images or audio-like signals. Log-frequency, log-spectrogram, and mel descriptors compress or reweight the spectrum and can reduce sensitivity to narrow high-frequency fluctuations while preserving broad-band energy trends. In the present grinding context, these descriptors should be treated as engineered spectral summaries rather than as psychoacoustic models, because the mel scale was not designed for AE physics. Their value is empirical and representational: they provide compact spectral summaries that may be useful when the number of grinding passes is limited.

Wavelet-scattering transforms offer another compromise between handcrafted features and learned deep representations. Scattering networks use cascaded wavelet modulus operators to generate translation-stable and deformation-stable descriptors without supervised learning of convolutional filters \cite{mallat2012group}, as demonstrated on various signal tasks \cite{Anden2014}. This property is useful when the dataset is too small to estimate many convolutional parameters reliably. In grinding, two-dimensional scattering descriptors computed from AE or vibration maps can provide structured multi-scale summaries while preserving more local information than scalar pass-level statistics. Their disadvantage is that the representation is fixed and may not match the most discriminative frequency bands for a particular machine, wheel, sensor, or material.

\subsection{Transparent and classical machine-learning models}\label{rw:transparent-ml}

Classical machine-learning approaches remain relevant for surface roughness prediction because many manufacturing datasets are small, structured, and tabular after feature extraction. Linear regression and shallow neural networks have been widely used for machining roughness prediction \cite{Benardos2003}, while support vector regression \cite{Ozel2005} and random forests or boosted trees \cite{Wuest2016} are preferred when process parameters and sensor summaries are available. These models are often easier to train and tune than deep networks and can be more stable when the number of experimental conditions is limited.

Ridge regression is a useful baseline because it regularizes correlated predictors and produces coefficients that can be inspected after standardization \cite{Hoerl1970}. A ridge model is not guaranteed to be physically correct, but it can reveal whether predictions rely on plausible process parameters or frequency bands. Random forests \cite{Breiman2001} and gradient-boosted trees \cite{Friedman2001} can model nonlinear interactions and threshold effects without requiring a large neural architecture. XGBoost \cite{Chen2016} and LightGBM \cite{Ke2017} provide efficient implementations of gradient-boosted decision trees and are frequently competitive on structured tabular tasks. Recent empirical comparisons have also shown that tree-based methods remain strong baselines for tabular datasets, especially when data are limited \cite{shwartz2022tabular} or heterogeneous \cite{grinsztajn2022tree}. This evidence supports the inclusion of tree ensembles in grinding roughness benchmarks, but it does not imply that they will always outperform deep models.

\subsection{Deep learning and multimodal fusion for manufacturing monitoring}\label{rw:deep-fusion}

Deep learning has been widely studied for manufacturing monitoring \cite{Wang2018}, machine fault diagnosis \cite{Khan2018}, and system health management \cite{Lei2020} because learned representations can reduce manual feature design. CNNs are suited to spectrograms and other two-dimensional maps; recurrent or temporal-convolutional models can process sequences; attention mechanisms and transformers can learn interactions across time steps, features, or modalities; and graph-based models can encode relationships among sensors, process variables, or operating conditions. These methods are especially appealing when the training set covers the expected operating range and when the test distribution resembles the training distribution.

Multimodal fusion is relevant because manufacturing sensors rarely provide independent or complete information. General sensor-fusion literature distinguishes data-level, feature-level, decision-level, and hybrid fusion strategies \cite{Khaleghi2013}. In grinding, feature-level fusion can combine AE descriptors, vibration descriptors, and process parameters before regression. Deep fusion architectures can learn modality-specific encoders before concatenation, bilinear interaction, attention pooling, or graph-based aggregation. Such models can capture cross-modal relations, but they also introduce additional parameters and design choices. When only several hundred passes are available, the gain from flexible fusion must be weighed against the increased risk of fold-specific tuning and unstable post-selection.

For this reason, deep-learning results in small manufacturing datasets should be interpreted as conditional on the experimental design, validation protocol, and tuning procedure. A ResNet-style spectrogram CNN may be appropriate when local time--frequency motifs generalize across conditions, whereas a pass-level tree ensemble may be preferable when the dominant information is contained in broad spectral summaries and process-regime variables. The present work therefore treats deep models as evaluated baselines, not as a class to be rejected. The comparison is intended to identify which representation--model combinations generalize under strict condition-level exclusion.

\subsection{Validation protocols, model-selection bias, and measurement uncertainty}\label{rw:validation}

Validation design is central in grinding roughness prediction because repeated passes within a condition share process settings and machine context. If such passes are split randomly, models may learn condition identity rather than condition-invariant roughness mechanisms. Similar issues have been discussed in cross-validation studies showing that model selection can produce biased performance estimates when tuning and evaluation are not separated \cite{varma2006bias}, leading to optimistic assessments \cite{Cawley2010}, and that structured data require validation folds aligned with the grouping structure \cite{Roberts2017}. Leave-one-condition-out validation is therefore more stringent than random K-fold validation for this manuscript because each test fold corresponds to a grinding condition unseen during training.

Condition-level validation also changes the interpretation of statistical testing. When the effective unit of generalization is the condition, fold-level errors are more relevant than pass-level errors for comparing models. This reduces the apparent sample size from hundreds of passes to 16 held-out conditions and makes strong statistical claims harder to justify. The present manuscript therefore distinguishes point-estimate ranking from inferential conclusions: the best point-estimate model is a random forest on fused AE and vibration dB-z descriptors, but the planned statistical comparison supports conclusive superiority only over the shallow MLP baseline.

Measurement uncertainty further constrains interpretation. If the best MAE is close to the manufacturer-quoted profilometer repeatability or below the expanded uncertainty scale, then very small differences between models may reflect metrology noise, trace placement, surface heterogeneity, or repeatability limits rather than meaningful predictive improvement. The absence of raw trace-level repeatability and gauge R\&R prevents a full measurement-system analysis. General measurement guidance therefore supports a cautious interpretation of sub-0.005--0.010~\(\upmu\)m ranking differences \cite{JCGM2008}. In this setting, robustness across conditions, failure-mode diagnosis, and physical plausibility may be more informative than small changes in aggregate MAE.

\subsection{Explainability for manufacturing machine learning}\label{rw:explainability}

Explainability in manufacturing monitoring is useful when models are intended to support process understanding, troubleshooting, or engineering review. Linear coefficients can show the direction and magnitude of standardized feature associations. TreeSHAP provides additive feature attributions for tree ensembles \cite{lundberg2017unified} and can summarize local and global dependencies \cite{Lundberg2020}. Grad-CAM can highlight regions of a spectrogram or time--frequency map that influence a CNN prediction \cite{selvaraju2017grad}. Broader XAI reviews emphasize that such methods improve model inspection but do not by themselves establish causal mechanisms \cite{Arrieta2020}.

This distinction is important in grinding. If SHAP values, ridge coefficients, or Grad-CAM maps emphasize AE bands, vibration bands, or process variables that are consistent with known contact or dynamic behaviour, the result should be described as a physically plausible association. It should not be described as proof that a specific grain-scale mechanism caused the measured roughness. The distinction is even more important for AE signals because sensor mounting, resonance, preamplifier response, and transfer paths can reshape the measured spectrum. Attributions above the nominal flat-response limit of the AE sensor should therefore be reported as model-relevant spectral regions requiring sensor-transfer-function verification rather than as direct evidence of high-frequency process physics.

\subsection{Uncertainty quantification under condition shift}\label{rw:uq}

Uncertainty quantification is increasingly used in machine learning to distinguish confident predictions from uncertain ones. MC-dropout provides a practical approximation to Bayesian predictive uncertainty by retaining dropout at inference time \cite{gal2016dropout}. Tree ensembles can provide empirical dispersion measures through out-of-bag predictions or tree-to-tree variance, although these quantities are not guaranteed to be calibrated predictive intervals. Broader reviews of uncertainty in deep learning show that uncertainty estimates can be useful but are often miscalibrated, particularly under distribution shift \cite{Abdar2021}. Dataset-shift studies similarly show that predictive uncertainty can degrade when test conditions differ from training conditions \cite{ovadia2019can}.

Distribution-free conformal prediction provides a formal route to finite-sample coverage under exchangeability assumptions, and extensions have been proposed for regression and certain shift settings \cite{Angelopoulos2023}. However, grouped grinding experiments challenge simple uncertainty calibration because only a small number of conditions are available and the hardest test condition may be unlike the calibration conditions. The present manuscript therefore treats uncertainty analysis as a diagnostic experiment. Full-LOGO MC-dropout intervals and RF OOB/tree-variance intervals under-cover difficult held-out conditions, especially Condition 7. This negative result is informative because it shows that accurate point prediction under average conditions does not imply reliable interval coverage under condition-level shift.

\subsection{Position of the present work}\label{rw:position}

The present study is positioned between physics-guided feature engineering and modern deep-learning fusion. It does not claim that deep learning is generally inferior for grinding monitoring. Instead, it tests whether more auditable physics-aware representations and classical models remain competitive when the evaluation excludes entire grinding conditions. This distinction matters because a random split can answer whether similar passes are predictable, whereas LOGO validation asks whether the learned mapping transfers to a new process regime.

Several aspects distinguish the benchmark. The dataset is a fractional 16-condition design with 319 valid passes, and every test fold is a previously unseen condition. The benchmark compares process parameters, time-domain statistics, handcrafted and physics-informed descriptors, dB spectrograms, dB-z descriptors, and two-dimensional wavelet-scattering transforms. Transparent or auditable baselines are evaluated together with shallow MLP, ResNet-style spectrogram CNNs, multimodal fusion CNNs, bilinear fusion, trajectory CNN, GNN fusion, TabTransformer, and attention-based regressors. Model explanations are interpreted as associations rather than causal proof, uncertainty intervals are evaluated as diagnostic outputs under grouped shift, and latency is restricted to feature-precomputed model inference rather than full edge deployment. This framing is intended to provide a cautious and reproducible comparison of representation--model choices for interpretable grinding roughness prediction.
